\section{An obstruction for arbitrary complexes on the graph supports}
\label{sec:localized-support}
We return to the norm-one graph supports of Proposition~\ref{prop:graph-obstruction}. Replacing that direct sum by extensions on the same supports does not remove its normal obstruction. The trace of multiplication by a normal coordinate detects this obstruction even when the individual support components are not assumed to deform.

\begin{theorem}\label{thm:localized-support}
Let $Z_1,\ldots,Z_r$ be translates of abelian subvarieties of the same codimension $c$ in a complex abelian variety $A$. Suppose every pairwise intersection has codimension at least two in each component. Let $E$ be a perfect complex whose cohomology sheaves are supported on $\bigcup_iZ_i$. Write
\[
 m_i(E)=\sum_j(-1)^j
 \operatorname{length}_{\mathcal O_{A,\eta_i}}
 \mathcal H^j(E)_{\eta_i}
\]
for its generic virtual multiplicity at $Z_i$. For a first-order deformation $\kappa\in H^1(A,T_A)$, let $\nu_i(\kappa)\in H^1(Z_i,N_{Z_i/A})$ be its normal restriction. If $m_i(E)\ne0$ and $\nu_i(\kappa)\ne0$ for some $i$, then $E$ has no perfect deformation in direction $\kappa$, even with an arbitrary infinitesimal gerbe twist.
\end{theorem}

\begin{proof}
Remove the pairwise intersections and then the other components. On the resulting open neighborhood of $Z_i^\circ$, the complex is supported on $Z_i^\circ$ alone. Since the removed subset has codimension at least two in the smooth variety $Z_i$, its local cohomology in degrees zero and one with coefficients in a locally free sheaf vanishes. The localization sequence gives an injection
\[
 H^1(Z_i,N_{Z_i/A})\lhook\joinrel\longrightarrow
 H^1(Z_i^\circ,N_{Z_i/A}|_{Z_i^\circ}).
\]
For finite intersections in a four-dimensional graph this is an isomorphism; only injectivity is needed.

Translate $Z_i$ to the origin and write $A=V/\Lambda$ and
$Z_i=W/\Lambda_i$, where $\Lambda_i=\Lambda\cap W$. Choose a complex
vector-space complement $N$ to $W$. A first-order deformation of a complex
torus is a first-order change of its periods. On a lifted coordinate cover
near $Z_i$, its transition maps can therefore be written
\[
 x_b=x_a+\lambda_{ab}+\epsilon a_{ab},\qquad
 z_b=z_a+\epsilon b_{ab},\qquad \epsilon^2=0.
\]
Here $x$ is tangential, $z$ is normal, $\lambda_{ab}\in\Lambda_i$, and
$a_{ab},b_{ab}$ are constant on each connected overlap. The cocycle
$(b_{ab})$ represents $\nu_i(\kappa)$. The tangential transitions are
independent of $z$ and define a first-order deformation of $Z_i$.
These coordinates may be taken on formal neighborhoods: the lattice of
the quotient torus is discrete, so only $\Lambda_i$ identifies points in
the completion along $Z_i$.

Suppose that a perfect lift $\widetilde E$ exists, allowing an arbitrary
gerbe whose central restriction is trivial. Trivialize the gerbe on the
coordinate charts and push the resulting complexes along the local
projection to the tangential chart. Denote the pushed complexes by $F_a$.
They are perfect. Indeed their cohomology is finite over the tangential
ring, because the normal ideal acts nilpotently on the central
cohomology and $\epsilon$ is nilpotent. A bounded complex of finite free
modules over the ambient formal power-series ring has finite Tor
amplitude over the tangential ring, since that power-series ring is flat.
Finite coherent cohomology and finite Tor amplitude imply perfection.
This argument permits arbitrary cohomology and does not require a flat
sheaf representative of $\widetilde E$.

The virtual rank of $F_a$ is $m=m_i(E)$. At the generic point of the
central tangential chart this is precisely the alternating length in the
statement, after forgetting the nilpotent normal-coordinate action.
For each normal coordinate take the categorical trace
\[
 t_a^\alpha=\operatorname{Tr}_{F_a}(z_a^\alpha).
\]
Modulo $\epsilon$ this trace is zero: over the generic point,
$z_a^\alpha$ acts nilpotently on every finite-dimensional cohomology
space, and its alternating trace vanishes. A holomorphic function on the
reduced chart vanishing generically vanishes identically. Thus
$t_a^\alpha=\epsilon f_a^\alpha$ for a holomorphic function
$f_a^\alpha$ on the central chart.

On an overlap, invariance of categorical trace under conjugation gives
\[
 t_b^\alpha-t_a^\alpha
   =\epsilon b_{ab}^\alpha\operatorname{Tr}_{F_a}(\operatorname{id})
   =\epsilon m b_{ab}^\alpha.
\]
The scalar gerbe factor on a triple overlap cancels under conjugation.
The tangential first-order translation changes $\epsilon f_a^\alpha$
only to order $\epsilon^2$. It follows that
$f_b^\alpha-f_a^\alpha=m b_{ab}^\alpha$. Hence
$m\nu_i(\kappa)|_{Z_i^\circ}=0$. The injection already proved and $m\ne0$
contradict the hypothesis.

The usual first-order obstruction theorem identifies this
nonexistence with the nonvanishing of
$\kappa\mathbin{\lrcorner}\operatorname{At}(E)+\beta\operatorname{id}_E$
for every $\beta$; see \cite[Corollary~3.4]{HT}. The scalar term is also
seen directly in the twisted \v Cech lifting equations. No simplicity
or vanishing of negative self-extensions is needed for this square-zero
obstruction.
\end{proof}

The local trace used here is consistent with the localized
semiregularity construction of Buchweitz--Flenner
\cite[Remark~4.7(2), p.~170; Proposition~8.2 and formula~(3),
pp.~200--201]{BF}. The generic multiplicity is the alternating length
specified in the theorem, not the virtual rank of the derived
restriction to the support, which has normal-Tor cancellation.

Distinct scalar norm-one graphs in $E_i^4\times E_i^4$ meet in finitely many points, and distinct translates of the same graph are disjoint. Their normal classes are $B^{\mathsf t}-B$. The theorem therefore excludes every perfect complex supported on a finite union of these graphs with nonzero generic virtual multiplicity from deforming in the six skew-matrix directions. In particular a nonzero pure degree-eight Weil character forces some such multiplicity to be nonzero. The restriction now covers arbitrary iterated cones, direct summands and thickenings on these supports. It leaves open complexes with larger support, different geometric supports, or components meeting along divisors.
